\subsection{Stable ordinary classes and the complete complex positional phase}

\begin{statement}{E.107}{Proved}{The fifth character, the exact exponent \(h=1\), and every complex toric fiber}
Use the literal roots, prime labels and ordinary class coordinates of
E.106. For the cubic \(B_{d+1}/B_d\), write
\(\chi_d:E_{B_d}\to\mathbb F_3^6\) for the six arithmetic local
norm characters, in the order \(Q_1,Q_2,Q_3,T_1,T_2,T_3\).
Every character uses the actual generator
\(\beta_{d+1}\mapsto\zeta\beta_{d+1}\). Put
\[
 S_Q(x)=x_1+x_2+x_3,\qquad S_T(x)=x_4+x_5+x_6,\qquad
 \mathcal H=\ker(S_Q+S_T).
 \tag{E107.1}
\]
The complete images are
\[
 \chi_2(E_{B_2})=\ker S_Q\cap\ker S_T,\qquad
 \chi_d(E_{B_d})=\mathcal H\quad(d\ge3).
 \tag{E107.2}
\]
Their dimensions are four and five, respectively. Five actual ordinary
units in \(B_3\), included in each later field, generate every latter
image. The proof constructs their classes from the complete E.106
image and an ordinary preimage of E.101's \(\eta\) under
\(N_{B_3/L_3}\). The latter unit has the exact triple sums
\((S_Q,S_T)=(1,2)\).

All later ordinary class groups and their actual norm maps stabilize:
\[
 \begin{gathered}
 \mathcal C_d=\operatorname{Cl}(B_d)\{3\}=C_9^2\oplus C_3
       \quad(d\ge3),\qquad
 D_d=\operatorname{Cl}(L_d)\{3\}=C_9a_d\oplus C_3b_d
       \quad(d\ge2),\\
 N_{B_{d+1}/B_d}=I_3,\quad j_{B_{d+1}/B_d}=3I_3
       \quad(d\ge3),\\
 N_{L_{d+1}/L_d}=I_2,\quad j_{L_{d+1}/L_d}=3I_2
       \quad(d\ge2).
 \end{gathered}
 \tag{E107.3}
\]
The coordinates are transported by the actual iterated ideal norms,
starting with E.106; these equalities specify the maps. The horizontal
ordinary class maps for every \(d\ge3\) are
\[
 N_{B_d/L_d}=
 \begin{pmatrix}3&6&6\\1&0&0\end{pmatrix},\qquad
 j_{B_d/L_d}=
 \begin{pmatrix}3&0\\3&3\\2&0\end{pmatrix}.
 \tag{E107.4}
\]
The \(h\)-action is the E.106 matrix
\(H=\left(\begin{smallmatrix}4&3&0\\8&7&6\\1&0&1\end{smallmatrix}\right)\);
the actual \(g_{d-1}\)-action is identity on \(\mathcal C_d\).
The six S-prime columns and the principal S-divisor lattice are
exactly \(C_S,H_S\) of E106.7--9 at every \(d\ge3\).
In particular
\[
 \operatorname{ord}[Q_i]=3,\qquad
 Q_1Q_2Q_3\ \hbox{is principal},\qquad
 q_d=[\mathfrak q_{3,d}]=3a_d,\quad
 \operatorname{ord}(q_d)=3\quad(d\ge2).
 \tag{E107.5}
\]
The actual horizontal ordinary unit norm is onto for every \(d\ge3\).
The vertical ordinary norm quotients are
\[
 E_{B_d}/N E_{B_{d+1}}\simeq C_3^6\quad(d\ge3),\qquad
 E_{L_d}/N E_{L_{d+1}}\simeq C_3^3\quad(d\ge2).
 \tag{E107.6}
\]
Their complete norm tests are the local field-norm characters followed,
on the zero character kernel, by the actual Hilbert--90 divisor quotient
\(\mathcal C_{d+1}/R_{d+1}=C_3\), or
\(D_{d+1}/\langle q_{d+1}\rangle=C_3^2\), respectively.
Every class coordinate is computed by the iterated norm to the fixed
E.106 class group. The completed norm-one kernels have types
\(C_3\oplus\mathbb Z_3^{2\cdot3^{d+1}}\) and
\(C_3\oplus\mathbb Z_3^{2\cdot3^d}\).

For the actual coherent three-unit group of E.69 one has
\[
 \begin{gathered}
 N_{L_d/K}\widehat{\mathcal O_{L_d}[1/3]^\times}_3
       =\lambda^{3\mathbb Z_3}\quad(d\ge2),\\
 I_0=\operatorname{pr}_0\mathcal U_\infty
       =\lambda^{3\mathbb Z_3},\qquad h=1 .
 \end{gathered}
 \tag{E107.7}
\]
Choose one actual coherent tuple
\(e\in\mathcal U_\infty\) with \(e_0=\lambda^3\), and put
\[
 \mathcal U_\infty^0=\ker(\operatorname{pr}_0:\mathcal U_\infty\to I_0).
 \qquad
 \mathbb Z_3\times\mathcal U_\infty^0
    \xrightarrow{\sim}\mathcal U_\infty,\quad(z,k)\longmapsto e^z k .
 \tag{E107.8}
\]
The inverse extracts \(z=v_\lambda(u_0)/3\) and then \(k=u e^{-z}\).
The compact kernel has the E.69 recursive coordinates
\[
                  \mathcal U_\infty^0
          \simeq\prod_{d\ge0}\mathbb Z_3^{\,2\cdot3^d}.
 \tag{E107.9}
\]
Changing \(e\) multiplies it by one element of this exact kernel.

The remaining complex positional phase is now fully evaluated.
Use \(\mathcal P,\mathscr R(\mathcal P),\mathcal K\) and the actual
classes \((\ell_p,m_p)\in C_9\oplus C_3\) of E.88, with
\(0\le\ell_p<9\). Define on its specified domain
\[
 \mathcal K=\left\{t\in\mathscr R(\mathcal P):
     \sum c_pt_p=\sum d_pt_p=0\pmod3\right\},\qquad
 \Theta_*(t)=-\frac{\sum_{p\in\mathcal P}\ell_pt_p}{3}.
 \tag{E107.10}
\]
Here \(c_p=\ell_p\bmod3\). The numerator belongs to \(3\mathbb Z_3\);
the sum converges in the restricted multiplication-by-\(3^m\)
topology. There is a continuous coherent divisor splitting \(s_*\)
with valuation zero at seventeen and phase exactly \(\Theta_*\).
The following are full forward and inverse coordinates:
\[
 \begin{aligned}
 \mathbb Z_3^2\times\mathcal U_\infty^0\times\mathcal K
      &\xrightarrow{\sim}\mathcal X,&
 (b,z,k,t)&\longmapsto\beta^b e^z k\,s_*(t),\\
 t&=\operatorname{div}_{\notin S}x,&
 b&=v_{17}(x_0),&
 z&=\frac{v_\lambda(x_0)-\Theta_*(t)}3,\\
 k&=x\beta^{-b}e^{-z}s_*(t)^{-1}.&&&
 \end{aligned}
 \tag{E107.11}
\]
In this formula \(v_{17}\) is the normalized valuation at the unique
prime of \(K\) above seventeen; \(\beta_0=17\zeta\).
The resulting actual rational base norm, in
\(\widehat{\mathbb Q^\times}_3\), is
\[
 N_{K/\mathbb Q}(x_0)
       =17^{2b}\,3^{\,3z-\frac13\sum\ell_pt_p}
                      \prod_{p\in\mathcal P}p^{t_p}.
 \tag{E107.12}
\]
Changing the normalized splitting while retaining the exact phase
\(\Theta_*\) is precisely multiplication by a continuous homomorphism
\(\mathcal K\to\mathcal U_\infty^0\). Both directions of this
description hold.

Here is the exact complete fiber statement in the section target \(H\).
Fix \(A=3\kappa_3(17)\) in E.46 and the pure positional component \(1\).
The image on the scaled toric coordinates consists exactly of
\[
 \begin{gathered}
 \left(\kappa_3(17),\
 B_{17}\kappa_3(17)+B_3\kappa_3(3)+
                  \sum_{p\in\mathcal P}r_p\kappa_3(p)\right),\\
 B_{17},B_3\in\mathbb Z_3,\qquad
 r\in\mathscr R(\mathcal P),\qquad
 \sum c_pr_p=0\pmod3.
 \end{gathered}
 \tag{E107.13}
\]
For every such pair its entire fiber is \(\mathcal U_\infty^0\).
The explicit map of that fiber and its inverse are
\[
 \begin{aligned}
 k\longmapsto x
   &=\beta^{\,3B_{17}/2}\,
     e^{\,B_3+\frac13\sum\ell_pr_p}\,
     k\,s_*(3r),\\
 k&=x\beta^{-3B_{17}/2}
     e^{-B_3-\frac13\sum\ell_pr_p}s_*(3r)^{-1}.
 \end{aligned}
 \tag{E107.14}
\]
Every exponent in this display is integral in \(\mathbb Z_3\).
The comparison to section classes is the actual Kummer--Shapiro
map E.46 and its total-norm cube condition; the pure component
has been specified. This evaluates the complete complex component
and its toric fibers, without claiming an intrinsic lift or an
evaluation of the separate pure-field norm system.
The manuscript anchor is LM \(\S51.16\), folios 733v--736v,
scans 172--178 through E.46, E.69, E.72 and E.88.
\end{statement}

\begin{proof}
\emph{Inclusion of consecutive actual character maps.}
For every \(d\ge2\), the extension
\(B_{d+1}/B_{d-1}\) is cyclic of degree nine. Its generator
\(g\) sends \(\beta_{d+1}\) to \(\xi\beta_{d+1}\);
its subgroup over \(B_d\) is generated by \(g^3\).
All six primes are totally ramified, with residue degree one
and the previously fixed unique vertical prolongations.
For \(u\in B_{d-1}^\times\) included in \(B_d\), the E.106
local Artin square consequently gives
\[
 \begin{aligned}
 3\chi_{d,w}(u)
  &=\chi_{9,v}(N_{B_{d,w}/B_{d-1,v}}u)
    =\chi_{9,v}(u^3)=3\chi_{9,v}(u)\pmod9,\\
 \chi_{9,v}(u)\pmod3&=\chi_{d-1,v}(u).
 \end{aligned}
 \tag{E107.15}
\]
The second equality is literal restriction: \(g\) sends
\(\beta_d=\beta_{d+1}^3\) to \(\zeta\beta_d\).
Division by the injection \(\mathbb Z/3\xrightarrow{\,3\,}\mathbb Z/9\)
therefore proves
\[
        \chi_d(u)=\chi_{d-1}(u)
              \quad\hbox{for every }u\in E_{B_{d-1}}.
 \tag{E107.16}
\]
Thus the images are nested, with precisely the same arithmetic
generators and labels. This is the actual norm/inclusion square,
obtained from the proved restriction/corestriction compatibility
of the Tate cup product
\cite[III,Theorem3.1 and its proof; III,\S4]{E_MilneCFT}.
It does not require an identification of different abstract norm
characters.

\emph{The actual fifth character and both split-prime sums.}
The complete matrix \(D\) in E106.3 has rank four and each of
its two triple sums is zero. Hence its image is exactly
\(\ker S_Q\cap\ker S_T\), an explicitly specified subspace
of dimension four.

Let \(\eta\in E_{L_2}\) be the actual ordinary unit constructed
in E101.19 and its following residue paragraph, and include it in \(L_3\). At seventeen its
residue lies in \(\mathbb F_{289}\) and satisfies
\[
                  \bar\eta^3=\zeta^2,\qquad
                  \bar\eta^{96}=\zeta .
 \tag{E107.17}
\]
The residue field and residue value remain unchanged under every
later totally ramified vertical inclusion. The tame unit
character at seventeen is calibrated in the established
arithmetic convention by \(\xi\): E106.3 gives
\(\chi_{2,T_i}(\xi)=1\), while
\(\bar\xi^{96}=\zeta^2\). Since the norm on residues of a
totally ramified tame cubic is cubing, its unit character is
therefore
\[
       \chi_{17}(a)=-\log_\zeta(\bar a^{96})\pmod3 .
 \tag{E107.18}
\]
This determines the character on the entire residue field, whose
cube quotient has order three. Principal units are cubes because
their pro-seventeen group admits unique division by three.
The uniformizer \(\beta_d\) has valuation one, so these are the
characters for the actual cubic radical generators at every
level. The residue calculation gives \(\chi_{17}(\eta)=2\).

For \(L_4/L_3\), all other finite primes are unramified and
\(\eta\) is an ordinary unit, so their Artin characters vanish.
All infinite places are complex. Global reciprocity with the
actual generator \(\beta_4\mapsto\zeta\beta_4\) thus gives
the character at the unique prime above three as \(1\).

E.106 proves the actual surjection
\(N_{B_3/L_3}E_{B_3}=E_{L_3}\). Choose \(v\in E_{B_3}\)
with \(N_{B_3/L_3}v=\eta\). Each S-prime of \(L_3\) splits
into three degree-one primes in the unramified extension
\(B_3/L_3\). The equality \(B_4=B_3L_4\) identifies the local
cubic groups by the same generator on \(\beta_4\).
Writing \(v_i\) for the three actual images in the common
completion, the local norm is \(\prod_i v_i\), and the
Artin maps are homomorphisms. Hence the two exact squares give
\[
 \begin{aligned}
 S_Q(\chi_3(v))
   &=\chi_{L_4/L_3,3}\left(\prod_{i=1}^3v_{Q_i}\right)
     =\chi_{L_4/L_3,3}(\eta)=1,\\
 S_T(\chi_3(v))
   &=\chi_{L_4/L_3,17}\left(\prod_{i=1}^3v_{T_i}\right)
     =\chi_{L_4/L_3,17}(\eta)=2 .
 \end{aligned}
 \tag{E107.19}
\]
There is no residue-degree multiplier in either square:
all three horizontal completions have degree one. The
horizontal automorphism fixes \(\zeta\) and \(\beta_4\), so
it does not change the target generator between the three maps.

This fifth vector is independent of the rank-four included image
in (E107.16). For every ordinary unit of \(B_d\), all local
characters away the six S-primes vanish. Global reciprocity
therefore gives \(S_Q+S_T=0\), proving the rank-five upper
bound at every stage. This proves (E107.2), including its
persistence under all later inclusions.
More explicitly, put \(w_0=(1,0,0,2,0,0)^t\).
The difference \(\chi_3(v)-w_0\) belongs to the actual rank-four
image. E.106's image inverse and ordinary unit preimages
therefore supply \(u\in E_{B_2}\) with that character.
Replacing \(v\) by \(v/u\) gives the literal fifth generator
of character \(w_0\). Together with four units realizing
the retained E.106 image columns, it proves the five-unit
assertion. The certificate records a coordinate inverse
for these five columns on the entire reciprocity hyperplane.

\emph{Stable ordinary classes and exact prime maps.}
The cyclic Hasse norm theorem converts the complete image
\(\mathcal H\) into ordinary-unit field-norm index \(3^5\)
for every \(B_{d+1}/B_d\), \(d\ge3\)
\cite[VIII,Theorem3.1 and proof]{E_MilneCFT}.
There are six ramified primes, each of ramification index three,
and no real places. E.91 proves that the actual vertical generator
acts identically on the entire ordinary three-class group.
The ambiguous class number formula therefore gives
\[
                  |\mathcal C_{d+1}|
                    =|\mathcal C_d|\frac{3^{6-1}}{3^5}
                    =|\mathcal C_d|\qquad(d\ge3).
 \tag{E107.20}
\]
The formula uses the field-norm index just evaluated
\cite[Theorem1(1) and its ideal/Herbrand derivation]{E_LemmermeyerAmbiguous}.

Every adjacent ordinary class norm is onto. Indeed, for either
ramified cyclic cubic \(E/F\), its intersection with the Hilbert
three-class field \(H_F\) is \(F\): otherwise the cyclic extension
would itself be unramified. In \(EH_F/F\), the Frobenius elements
\((1,a)\) give degree-one primes of \(E\) whose ideal norms
represent every \(a\in\operatorname{Gal}(H_F/F)\).
This is the actual ideal-norm/Artin square, not an order comparison
without a map
\cite[V,Proposition3.3,Theorem3.6,Example3.9;
VIII,Theorem7.4]{E_MilneCFT}.
Consequently (E107.20) makes all B norms from level three
isomorphisms.

The unique vertical S-prime prolongations have residue degree one,
so these isomorphisms send each named prime to the same named
prime at the preceding level. Every prime column in E106.7 is
therefore retained, and its integral relation lattice is exactly
E106.8--9 at every level. All these prime classes are three-primary,
as proved there from their powers extending B1 primes.
Thus these relations are actual principality. In particular every
\(Q_i\) has order three and their product is principal.
The horizontal norm sends \(Q_i\) to \(\mathfrak q_{3,d}\);
its norm down to \(L_2\) has order three. The original prime
therefore has order exactly three at every \(d\ge2\).
E91.3 now gives \(D_d=C_9\oplus C_3\). Its surjective norms
are isomorphisms too.

Transporting the actual coordinate bases by these norm
isomorphisms proves every identity in (E107.3).
In particular the vertical actions on \(D_d\) are identity,
and the ideal identity \(jN=\sum\sigma^i\) gives \(j=3I\).
On the B groups this action identity was already E.91.
The restriction of \(g_{d-1}\) to \(B_{d-1}\) is the vertical
generator there, whose class action is identity. Norm equivariance
and its injectivity prove that \(g_{d-1}\) also acts identically
for \(d\ge4\); the case \(d=3\) is the vanishing \(W=0\) in E.106.
The \(h\)-action transports to the same matrix \(H\).

The norm squares commute by transitivity. Extension also commutes
with the other vertical norm in the horizontal square: in this
linearly disjoint square, \(h\) and the vertical generator commute,
and the prime ramification/residue coefficients agree; equivalently
it is the ideal restriction/transfer identity for the two subgroups.
Thus both horizontal matrices remain E106.5. Their extension
kernel has order three at every level. Hilbert--90 for principal
ideals in this unramified cubic, and the unit Herbrand quotient
\(1/3\), give
\[
 3=|\ker j_{B_d/L_d}|
   =3[E_{L_d}:N_{B_d/L_d}E_{B_d}]\quad(d\ge3).
 \tag{E107.21}
\]
This proves all the horizontal ordinary norm surjections.
The Herbrand quotient here is the ordinary-unit case of the
proved lattice comparison
\cite[VII,Proposition3.1 and its complete proof]{E_MilneCFT}.

The six S-primes generate the same order-\(81\) subgroup
\(R_d\subset\mathcal C_d\). The strongly ambiguous subgroup
for an adjacent B step is \(j\mathcal C_d+R_{d+1}=R_{d+1}\);
the inclusion \(3\mathcal C_{d+1}\subset R_{d+1}\) follows
from the displayed prime columns. Its quotient in the full
ambiguous group has order three and is the coordinate \(y\bmod3\).
The exact divisor Hilbert--90 map of E106.20 therefore gives the
additional ordinary-unit norm index three. Multiplied by the
field index \(3^5\), this is \(3^6\).

For an L step, the included \(\eta\) has noncube residue at seventeen.
Global reciprocity at the two ramified primes bounds its
ordinary-unit field-norm image by dimension one and this residue
makes that dimension exactly one. The strongly ambiguous subgroup
is \(3D_{d+1}=\langle q_{d+1}\rangle\), of order three;
the other ramified prime is the principal divisor of \(\beta_{d+1}\).
The remaining quotient has order nine. The full ordinary
norm index is therefore \(3\cdot9=27\). These quotient groups
are killed by three because an included base unit has norm its
cube, proving (E107.6).
The same explicit Hilbert--90 divisor argument as E.106 proves
necessity and sufficiency of each stated two-step test.
Completion preserves these finite-index norm images.
Unit ranks and the cubing kernels on \(\mu_9\) or \(\mu_3\)
give the stated completed kernel types.

\emph{The exact coherent three-unit exponent.}
The distinguished prime now has order three at every \(d\ge2\).
The principal-valuation/norm identity of E.72 consequently gives
the first equality of (E107.7). There is no missing root-of-unity
coordinate: all these images factor through the literal
\(N_{L_1/K}\mathcal U_1=\lambda^{\mathbb Z_3}\) of E69.16.

For clarity, this finite-level assertion supplies a compatible
tuple of base norm \(\lambda^3\). Each completed three-unit group
\(\mathcal U_d\) is compact. In their product impose the closed
adjacent norm equations and the equation \(u_0=\lambda^3\).
Any finite set of these conditions is satisfied by an element
at the largest specified level, of norm \(\lambda^3\), and all
its lower norms. The finite-intersection property gives a tuple
in the full product. Conversely every coherent tuple projects
to each finite norm image, hence to \(\lambda^{3\mathbb Z_3}\).
This proves \(I_0=\lambda^{3\mathbb Z_3}\) and \(h=1\) with
the actual transition maps.
It proves the existence of \(e\) in (E107.8).
The displayed maps there are continuous inverses; changing \(e\)
uses exactly \(\ker\operatorname{pr}_0\).
The already proved recursive norm-kernel construction E69.6
gives (E107.9), with its specified continuous forward and inverse
maps. No compactness is asserted for a restricted divisor module.

\emph{Exact phase normalization and every toric fiber.}
Start with an actual continuous splitting \(s\) supplied by
E88.7, normalized to have valuation zero at seventeen.
Its phase \(\Theta(t)\) satisfies
\(\Theta(t)\equiv-\sum\ell_pt_p/3\pmod3\) by E88.10.
The right side is an integral continuous linear map on
\(\mathcal K\): its numerator is divisible by three because
\(\ell_p\equiv c_p\pmod3\), and the restricted sequence
condition makes its sum convergent. Define
\[
  q(t)=\frac{\Theta_*(t)-\Theta(t)}3\in\mathbb Z_3,\qquad
                s_*(t)=s(t)e^{q(t)} .
 \tag{E107.22}
\]
This is an actual continuous coherent map. It preserves the
outside-S divisor and the seventeen valuation, and adds \(3q(t)\)
to the base three valuation. Hence its phase is exactly \(\Theta_*\).
Two splittings with these same data differ by a continuous
homomorphism to \(\mathcal U_\infty\) with base valuation zero.
Since its base image is exactly \(\lambda^{3\mathbb Z_3}\),
that condition is precisely membership in \(\mathcal U_\infty^0\).
Conversely multiplying by any such homomorphism preserves all
the data. This proves the full normalization ambiguity.

Now E88.7 and (E107.8) give (E107.11), including its inverse:
after removing the outside divisor and the seventeen coordinate,
the remainder is a coherent three-unit, whose base valuation is
divisible by three. The compact factors are continuous under
\(\mathbb Z_3\)-powers. The restricted factor is continuous by
the convergent products proved in E.88, and (E107.22) is a
continuous correction. These facts prove the asserted
topological isomorphism and all inverse-coordinate claims.
Taking the actual norm of the zeroth coordinate uses
\(N\beta_0=17^2\), \(N\lambda=3\), and degree-one norm \(p\)
for each chosen outside-S prime, giving (E107.12).
Prime-to-three rational torsion has already disappeared in
\(\widehat{\mathbb Q^\times}_3\).

The total-norm cube condition in E.46 can now be tested
coefficient by coefficient in that same completed multiplicative
group. Equation (E107.12) is a cube exactly when
\[
 b\in3\mathbb Z_3,\qquad t=3r,\qquad
 \sum\ell_pr_p=0\pmod3
       \quad\Longleftrightarrow\quad \sum c_pr_p=0\pmod3 .
 \tag{E107.23}
\]
When \(t=3r\), both defining class congruences for \(t\in\mathcal K\)
already hold. Dividing the norm exponents by three, as prescribed
by the actual toric map E46.7, gives
\[
       B_{17}=2b/3,\qquad
       B_3=z-\frac13\sum\ell_pr_p .
 \tag{E107.24}
\]
These two equations have the unique inverse used in (E107.14),
because two is a unit in \(\mathbb Z_3\).
All outside coefficients are the specified \(r_p\).
The remaining coordinate is exactly \(k\in\mathcal U_\infty^0\),
with no additional stabilizer in the E.46 section comparison.
This proves both directions of (E107.13) and the complete
fiber statement with its actual maps.

The certificate exhausts the rank-four kernel, all eighty-one
possible fifth vectors with triple sums \((1,2)\), the entire
reciprocity hyperplane and its coordinate inverse. It checks
all residue units and their character products in
\(\mathbb F_{289}\), every stable B class, the principal S-lattice,
the full retained seventy-nine-prime phase basis, and the
complete modulo-nine cube condition on five actual prime
coordinates. The all-level conclusions follow from the norm
squares and induction just proved; no upper BNF is used.
\end{proof}
